\section{工程验收与回归体系：reference、测试、覆盖率、性能门禁和框架对标}

\subsection{为什么“能跑”不是完成}

一台本地大模型推理引擎最容易给人错觉：只要 prompt 能生成文字，项目就算跑通。对学习项目来说，这只是第一步；对高质量 AI infra 来说，这远远不够。生成文本看起来正常，不代表 tokenizer 合同正确，不代表量化 scale 正确，不代表 KV cache 没有越界，不代表 SIMD tail path 被覆盖，不代表 GPU fallback 没有偷偷发生，也不代表性能优化真的击中了瓶颈。

工程验收要回答四个问题：

\begin{enumerate}
  \item 正确性由哪个 reference 定义。
  \item 哪些层次的 oracle 已经比较过。
  \item 哪些性能指标和硬件环境可复现。
  \item 哪些回归测试会在后续改动时保护这些事实。
\end{enumerate}

\begin{keyidea}
AI infra 的完成标准不是“输出了 token”，而是 reference、oracle、profiler、覆盖率、性能回归和外部框架对标形成闭环。没有闭环，优化无法被信任，也无法长期维护。
\end{keyidea}

\subsection{reference 分层：先有正确答案}

reference 是不追求速度的正确性来源。它应该清晰、可读、可调试。后端优化越激进，reference 越重要。项目中至少需要下面几类 reference：

\begin{center}
\begin{tabular}{p{0.22\linewidth}p{0.34\linewidth}p{0.30\linewidth}}
\toprule
reference 层次 & 作用 & 不追求的东西 \\
\midrule
字节/量化 converter & 定义 int4/int8 bytes 如何还原 & SIMD 和吞吐 \\
单算子 reference & 定义 RMSNorm/linear/attention 输出 & cache 和线程效率 \\
层级 reference & 定义 decoder block 输出 & 融合和内存复用 \\
端到端 reference & 定义 logits/生成序列合同 & 生产级速度 \\
\bottomrule
\end{tabular}
\end{center}

reference 代码要尽量少依赖后端复杂路径。比如 int4 reference converter 不应该复用 AVX2 解包函数；否则 SIMD 解包错了，reference 也跟着错。单算子 reference 可以用标量循环和 fp32 累加；端到端 reference 可以只跑小模型、小层数或模型切片，保证测试可承受。

\subsection{oracle 分层：从字节到生成序列}

oracle 消费 reference 和被测实现，判断误差是否在合同内。不同层 oracle 关注不同问题：

\begin{lstlisting}[numbers=none,caption={oracle 层次与定位能力}]
byte oracle:
  packed bytes -> decoded values
  catches nibble order, zero point, scale, tail

operator oracle:
  optimized kernel output vs scalar reference
  catches SIMD, accumulation, layout, quant

layer oracle:
  decoder block output vs reference block
  catches fusion, arena overwrite, state edge

logits oracle:
  final logits distribution
  catches visible model quality drift

sequence oracle:
  fixed sampler and seed token sequence
  catches end-to-end behavior changes
\end{lstlisting}

误差容忍度必须和 dtype、量化、后端和 reduction 顺序匹配。fp32 reference 对 fp32 scalar 可以要求很严；fp16、bf16、int8、int4、KV cache 量化都需要不同 tolerance profile。oracle 报告里要写清使用的是哪个 profile，不要把阈值埋在测试代码里。

\subsection{测试矩阵：不要只测漂亮形状}

AI 后端最容易漏测的是边界形状。测试矩阵要覆盖整齐形状，也要覆盖 tail、非对齐、不同 group size、不同 context length 和不同阶段。

\begin{lstlisting}[numbers=none,caption={后端测试矩阵示例}]
shapes:
  H = 4096
  H = 4097 for tail test
  O not divisible by tile_n
  prompt length: 1, 7, 32, 512
  context length: 1, 128, 1024, 4096

dtypes:
  f32 reference
  f16/bf16 style
  int8 per-channel
  int4 group-size 32/64/128

layouts:
  contiguous
  padded
  packed CPU panel
  device layout

backends:
  scalar
  CPU SIMD
  CPU threaded
  GPU kernel
  fallback path
\end{lstlisting}

这些测试不一定全部在每次提交跑。可以分层：快速单元测试常跑，较大的端到端和性能测试在夜间或发布前跑。但测试矩阵必须被设计出来，否则项目只会在“常见 happy path”上看起来稳定。

\subsection{coverage：覆盖率要服务风险，不服务数字}

用户要求每一步优化都做好，C++20 项目也应该有覆盖率意识。但覆盖率不是盲目追数字。高风险路径要优先覆盖：verifier 错误诊断、descriptor 解析、quant converter、tail path、arena liveness、KV cache position、fallback、sampler 和 profiler report。

\begin{center}
\begin{tabular}{p{0.22\linewidth}p{0.34\linewidth}p{0.30\linewidth}}
\toprule
模块 & 高风险分支 & 覆盖重点 \\
\midrule
model/verifier & 坏 shape、坏 dtype、坏 quant & 拒绝输入和诊断质量 \\
quant & tail、scale shape、packed order & reference converter 和 oracle \\
memory planner & last use、alignment、debug materialization & offset 不重叠和生命周期 \\
backend CPU & tail、非对齐、多线程分区 & 正确性和性能回归 \\
backend GPU & fallback、staging、sync & 端到端 latency 和资源释放 \\
runtime & prefill/decode 切换、stop、context limit & session 状态不串扰 \\
\bottomrule
\end{tabular}
\end{center}

如果项目工具链支持覆盖率，核心模块应该接近或超过 90\% 的有意义覆盖。若某些 GPU 分支或硬件计数器难以在 CI 中覆盖，要用 fake backend、mock event、small fixture 或本地验收脚本补充。重要的是把不可覆盖原因写清，而不是假装没有风险。

\subsection{CI 分层：不是所有测试都每次全跑}

AI 推理项目的测试成本差异很大。字节级 quant converter 可以毫秒级完成；7B 模型长上下文 benchmark 可能需要专门机器。工程上要分层，而不是把所有检查都塞进每次提交，也不能因为全量测试太慢就放弃测试。

\begin{center}
\begin{tabular}{p{0.12\linewidth}p{0.30\linewidth}p{0.38\linewidth}}
\toprule
层级 & 触发时机 & 内容 \\
\midrule
L0 & 本地开发频繁运行 & 格式化、编译、descriptor/verifier 单测、字节级 quant 测试 \\
L1 & 每次提交或 PR & 小模型 reference、单算子 oracle、memory planner、CPU scalar/SIMD 小矩阵 \\
L2 & 夜间或合并前 & 模型切片端到端、长 context、paged KV、fallback、覆盖率报告 \\
L3 & 发布前或专用机器 & 真实模型 benchmark、GPU、框架对标、p50/p95、峰值内存、质量回归 \\
\bottomrule
\end{tabular}
\end{center}

失败报告也要结构化。不要只贴一段日志。最小字段可以是：

\begin{lstlisting}[numbers=none,caption={CI failure 结构}]
CiFailure:
  layer: L0 | L1 | L2 | L3
  component: verifier | quant | backend | runtime | benchmark
  fixture_id
  backend
  expected
  actual
  tolerance_or_gate
  first_bad_artifact
  reproduction_command
\end{lstlisting}

这样开发者能直接知道问题属于哪层合同。若 byte oracle 失败，先查 nibble order、scale 和 tail；若 layer oracle 失败，查 fusion、arena 和状态边；若 L3 benchmark 回归，先看硬件、热身、频率、prompt set 和 profiler breakdown。CI 的作用不是制造红叉，而是把问题定位到正确工程层。

\subsection{性能门禁：基线、阈值和报告要稳定}

性能回归比功能回归更难测，因为机器负载、频率、温度和系统噪声都会影响数字。工程上不能用单次最好成绩做门禁，而要固定环境、固定输入、固定 warmup 和迭代，并报告分位数。

\begin{lstlisting}[numbers=none,caption={性能门禁报告字段}]
PerformanceGate:
  hardware_id
  backend_plan_id
  model_fixture_id
  quant_profile_id
  prompt_set_id
  thread_count_or_device
  warmup_iterations
  measurement_iterations
  baseline_version
  current_version
  allowed_regression_percent
  metrics:
    time_to_first_token
    prefill_tokens_per_s
    decode_tokens_per_s
    p50_step_latency
    p95_step_latency
    peak_memory
\end{lstlisting}

门禁阈值要分指标。比如 peak memory 回归 20\% 可能不能接受；单次 p50 latency 回归 3\% 可能在噪声内；p95 latency 回归 15\% 需要调查；pack time 变慢但 decode 变快是否可接受，要看目标场景是短请求还是长期服务。性能门禁不是机械判分，而是让退化可见。

\subsection{服务化 SLO：把 kernel 数字换成线上账本}

底层 kernel 快，不等于服务能扛业务。推理服务还要处理排队、取消、流式返回、上下文容量、KV cache 内存和 OOM。最小 serving 验收要把单请求 trace 聚合成下面的账本：

\begin{lstlisting}[numbers=none,caption={ServingSloEnvelope}]
ServingSloEnvelope:
  max_waiting_requests
  max_running_requests
  admission_policy: reject | queue | degrade
  ttft_p50_ms
  ttft_p95_ms
  decode_step_p95_ms
  end_to_end_p95_ms
  queue_wait_p95_ms
  kv_cache_reserved_bytes
  kv_cache_used_bytes
  oom_reject_count
  cancel_reclaim_p95_ms
\end{lstlisting}

用 7B 典型 shape 粗算一次 KV 容量：\code{layers=32, kv\_heads=32, head\_dim=128, fp16} 时，每 token KV 字节为：

\[
2 \times 32 \times 32 \times 128 \times 2 = 524288 \text{ bytes} = 512 \text{ KiB}
\]

一个 4K context session 约 \code{2 GiB} KV cache。若机器只给 serving 预留 \code{16 GiB} KV 空间，理论上最多同时容纳 8 个满长 session；实际还要扣除 allocator 碎片、paged metadata、workspace 和安全水位。admission control 因此不能只看请求数，必须看 \code{prompt\_tokens + max\_new\_tokens} 对 KV budget 的占用。

调度也要定量。若单请求 decode step p95 是 \code{40 ms}，同一 worker 串行跑 4 个活跃 session，最坏情况下某个 session 下一个 token 要等另外 3 个 step，队列等待可接近 \code{120 ms}。若业务要求流式 token 间隔 p95 小于 \code{80 ms}，就必须降低同 worker 并发、改 continuous batching、或把长请求降级/拒绝。这里的重点不是精确排队论，而是 SLO 要能反推出调度策略。

\begin{center}
\begin{tabular}{p{0.22\linewidth}p{0.34\linewidth}p{0.28\linewidth}}
\toprule
线上现象 & 应记录的字段 & 常见根因 \\
\midrule
TTFT 高 & queue wait、load、prefill、pack time & 冷启动、长 prompt、排队 \\
token 间隔抖动 & decode p95、active sessions、batch size & 调度过载、同步点、抢占 \\
OOM 或频繁拒绝 & KV reserved/used、workspace、arena peak & admission 不看 token 预算 \\
取消后内存不降 & cancel reclaim、session state、KV refcount & 资源生命周期错 \\
p50 正常但 p95 差 & queue wait、page fault、fallback count & 尾部排队、缺页、后端降级 \\
\bottomrule
\end{tabular}
\end{center}

因此服务化验收至少要有三条故障演练：超长上下文被 admission 拒绝并返回 token budget；取消请求后 KV、arena、workspace 在限定时间内释放；强制注入慢请求后，报告能拆出 queue wait 与 kernel time，而不是只看到端到端超时。

\subsection{profiler schema：优化证据要可比较}

profiler 输出不能每章一个格式。统一 schema 让不同优化可以比较。最小 schema 应该覆盖 stage、segment、backend、context、耗时、内存和可选硬件计数器。

\begin{lstlisting}[numbers=none,caption={统一 profiler event schema}]
ProfilerEvent:
  run_id
  stage: load | prepare | prefill | decode | sample
  segment_id
  layer_id
  backend
  kernel_or_operation
  start_ns
  end_ns
  input_tokens
  context_length
  bytes_read_estimate
  bytes_written_estimate
  arena_peak_bytes
  workspace_bytes
  kv_cache_bytes
  counters_optional
\end{lstlisting}

同一 schema 可以服务 CPU 和 GPU。CPU 可能填硬件计数器；GPU 可能填 stream event 时间和 staging bytes；reference 路径也可以填 stage 和耗时。这样后续报告不需要人工拼日志。

\subsection{诊断质量：错误要能定位到合同}

高质量工程不是永远不出错，而是出错时能定位。诊断应该回到合同，而不是只报底层异常。

\begin{lstlisting}[numbers=none,caption={诊断应该包含的信息}]
Diagnostic:
  severity
  component: manifest | verifier | compiler | backend | runtime
  object: tensor/operator/segment/session
  expected
  actual
  derived_from
  suggested_next_check
\end{lstlisting}

例如 backend 不支持某个 q4 格式时，应该写明 tensor、format、group size、backend capability 和 fallback 选择；KV cache 越界时，应该写明 session position、capacity、max context 和触发的 token；arena 复用冲突时，应该写明两个 value 的生命周期和 offset。这样的诊断能让读者把错误和书中讲过的合同对应起来。

\subsection{框架对标：尊重成熟实现，也要知道差距在哪里}

自研引擎需要对标成熟框架，但对标不能只贴一个 tokens/s。成熟框架可能使用不同量化格式、不同 tokenizer、不同线程数、不同 GPU kernel、不同 batch 策略。对标前要固定变量：

\begin{itemize}
  \item 同一模型或同一模型切片。
  \item 同一 tokenizer 和 prompt set。
  \item 同一量化格式或明确说明格式差异。
  \item 同一线程数、亲和策略或同一 GPU。
  \item 同样区分加载、prepare、prefill、decode 和 sampler。
  \item 同样报告正确性或至少 logits/生成差异。
\end{itemize}

对标表应该能说明差距：

\begin{center}
\begin{tabular}{p{0.20\linewidth}p{0.24\linewidth}p{0.24\linewidth}p{0.20\linewidth}}
\toprule
指标 & 本项目 & 成熟框架 & 解释 \\
\midrule
load/prepare & A & B & mmap、packing、上传差异 \\
prefill & A & B & GEMM/attention kernel 差异 \\
decode & A & B & GEMV、KV cache、launch 差异 \\
memory peak & A & B & arena、packed weight、KV 策略 \\
quality drift & A & B & 量化和 sampler 差异 \\
\bottomrule
\end{tabular}
\end{center}

如果差距很大，也不是失败。高质量书稿应该教读者定位差距：是算法、layout、kernel、线程、设备同步、量化、还是 runtime 调度。只要能定位，就能进入下一轮优化。

\subsection{发布前验收清单}

第三册对应的 C++20 项目在发布一个阶段性版本前，至少应该完成下面清单：

\begin{rubricbox}
\begin{itemize}
  \item 模型 manifest、tokenizer 合同、shape/dtype/quant verifier 有负例测试。
  \item reference converter、单算子 reference、层级 reference 和小端到端 reference 可运行。
  \item CPU scalar、CPU optimized、GPU 或 fallback 后端都通过 oracle。
  \item activation arena、workspace、packed weight、KV cache 有生命周期和容量测试。
  \item prefill/decode 分别有 benchmark，报告 time to first token、tokens/s、latency 分位数和内存峰值。
  \item 量化报告同时包含容量、scale metadata、logits/top-k、生成序列和性能。
  \item profiler event schema 稳定，报告能追溯到 segment、kernel、backend 和 context length。
  \item 性能回归门禁记录硬件、线程、prompt、baseline 和允许退化阈值。
  \item 与成熟框架的对标至少能解释最大差距来源。
\end{itemize}
\end{rubricbox}

这份清单看起来严格，但它正是 AI infra 工程和普通 demo 的区别。demo 证明“能跑”；验收体系证明“能长期改、能定位错、能解释快慢、能对用户负责”。

\subsection{阶段发布报告模板}

每个阶段发布时都应该留一份短报告。模板可以固定，内容不需要冗长，但必须能复现。

\begin{lstlisting}[numbers=none,caption={阶段发布报告模板}]
ReleaseReport:
  version
  model_scope
  supported_backends
  supported_quant_profiles
  correctness:
    reference_commit
    oracle_profiles
    failed_or_skipped_cases
  performance:
    hardware
    prompt_set
    prefill_tokens_per_s
    ttft_ms
    decode_tokens_per_s
    p95_step_latency_ms
    memory_peak_bytes
  known_limits:
    max_context
    unsupported_dtypes
    unsupported_shapes
    fallback_paths
  comparison:
    baseline_version
    mature_framework_reference
    largest_gap_explanation
\end{lstlisting}

这个报告能直接回答三个问题：能跑哪些模型和 profile，正确性证据是什么，性能瓶颈还在哪里。若报告写不出来，说明项目某些合同还没有建立好。

\subsection{验收失败时的处理顺序}

失败后不要直接调阈值。处理顺序应该是：

\begin{enumerate}
  \item 先确认 fixture、模型、prompt、backend plan 和量化 profile 是否一致。
  \item 再确认 reference 是否可信，尤其 byte converter 和 tokenizer。
  \item 定位最早失败层级：byte、operator、layer、logits、sequence、performance。
  \item 若是 correctness，先修实现或合同；若是性能，先确认测量口径和硬件状态。
  \item 只有在误差来源被解释后，才调整 tolerance profile。
\end{enumerate}

阈值是合同，不是让测试通过的旋钮。尤其是量化和 GPU fallback，不能因为输出文本看起来还行就放宽 logits/top-k 检查。

\subsection{一次 CPU kernel 优化的发布门禁示例}

假设本周把 \code{linear} 的 CPU 路径从标量 baseline 换成 packed int8 SIMD kernel。一个合格发布门禁不应该只跑一个 benchmark，而要按证据链通过。

\begin{lstlisting}[numbers=none,caption={CPU int8 kernel 发布门禁}]
1. byte oracle:
   int8 weights and scale metadata decode to expected values

2. packed weight oracle:
   packed panel maps back to original W[out, in]

3. operator oracle:
   scalar reference vs packed SIMD output

4. layer oracle:
   decoder block output under fixed fixture

5. logits oracle:
   final logits and top-k overlap

6. performance gate:
   decode latency, p95, memory peak and pack time

7. fallback check:
   unsupported ISA falls back to scalar or refuses at plan time
\end{lstlisting}

这条门禁能定位失败。例如 byte oracle 失败，说明量化字节解释有错；packed weight oracle 失败，说明 layout 转换有错；operator oracle 失败，说明 SIMD lane、accumulator、tail 或 scale 有错；layer oracle 失败，可能是 arena 复用、fusion 或状态边；performance gate 失败，才进入 cache、线程、NUMA、带宽和反汇编分析。

\subsection{一次 KV cache bug 的定位示例}

KV cache bug 很常见，因为它跨 token 存活，且 shape 往往看起来正确。假设长上下文续写时 logits 从第 200 个 token 开始漂移，可以按下面顺序定位：

\begin{lstlisting}[numbers=none,caption={KV cache 漂移定位流程}]
check tokenizer:
  token ids match reference

check model position:
  each token position is monotonic and not reset by cache slot

check RoPE checkpoint:
  q/k after rope match reference at positions 0, 1, 127, 200

check KV append:
  key/value written at expected logical position

check KV read:
  query head maps to correct kv head
  attention reads positions 0..current_position

check attention probability:
  score and softmax drift starts at which position
\end{lstlisting}

如果只比较最终生成文本，很难知道问题在哪。把 checkpoint 放在 RoPE、KV append、KV read 和 attention probability，能把“长文本质量差”拆成具体合同错误。

\subsection{一次性能回归的排查示例}

性能回归也要按阶段拆。假设版本 B 比版本 A 的 decode tokens/s 下降 18\%，不要直接改 kernel。先确认报告坐标一致：

\begin{itemize}
  \item 模型、prompt set、quant profile、backend plan 是否相同。
  \item 线程数、亲和、频率、温度和系统负载是否可比。
  \item pack time 是否被算进 decode。
  \item fallback count 是否增加。
  \item context length 分布是否变化。
  \item p50 和 p95 是否一起下降，还是只有尾延迟变差。
\end{itemize}

随后看 profiler breakdown：

\begin{center}
\begin{tabular}{p{0.24\linewidth}p{0.34\linewidth}p{0.26\linewidth}}
\toprule
现象 & 可能原因 & 下一步证据 \\
\midrule
权重 segment 变慢 & packing layout、SIMD path、NUMA & packed oracle、perf counters \\
KV segment 随 context 变差 & KV layout、cache miss、TLB & context bucket 报告 \\
sampler 时间上升 & logits copy、top-k 实现 & sampler trace \\
p95 上升但 p50 稳定 & 调度、队列、页错误 & queue wait、page fault \\
GPU kernel 快但端到端慢 & copy、sync、launch、fallback & stream event trace \\
\bottomrule
\end{tabular}
\end{center}

这类排查流程应该写进 release report 的失败处理说明。否则团队会在每次回归时重新发明定位方法。

\subsection{验收报告中的最小可读示例}

报告不需要长，但要能回答“这次改动可信吗”。下面是一个压缩示例：

\begin{lstlisting}[numbers=none,caption={release gate 摘要示例}]
ReleaseGateSummary:
  change: cpu int8 packed linear kernel
  model_fixture: tiny_decoder_v3
  backend_plan: cpu_avx2_int8_pack_v1
  correctness:
    byte_oracle: pass
    packed_weight_oracle: pass
    operator_oracle_max_abs: 0.0061
    logits_topk_overlap: 9/10
  performance:
    decode_tokens_per_s: 18.4 -> 24.7
    p95_step_latency_ms: 68.2 -> 49.5
    pack_time_ms: 143
    memory_peak_mb: 812 -> 856
  risk:
    only group size 64 covered
    avx512 path not implemented
\end{lstlisting}

这样的摘要比“快了 34\%”更有价值。它同时暴露了收益、误差、内存代价和限制。

\subsection{本节验收线}

读完本节，应该能建立工程闭环：

\begin{itemize}
  \item reference 定义正确答案，oracle 分层比较误差，profiler 证明性能变化。
  \item 测试矩阵必须覆盖 tail、非对齐、量化 group、context length、fallback 和不同后端。
  \item 覆盖率要优先服务 verifier、quant、memory planner、backend、runtime 等高风险路径。
  \item 性能门禁要固定硬件、输入、warmup、迭代、baseline 和允许退化阈值。
  \item profiler schema 要统一，让 CPU/GPU/reference/fallback 报告可以比较。
  \item 诊断要回到 manifest、IR、descriptor、backend plan 和 session 状态这些合同。
  \item 框架对标要拆开 load、prepare、prefill、decode、memory 和 quality，而不是只报一个总分。
\end{itemize}

至此，第六章形成闭环：CPU 和 GPU 后端不是孤立优化技巧，而是以 C++20 资源边界为基础，以编译器 lowering 为入口，以 oracle/profiler/回归门禁为出口的工程系统。
